\documentclass[11pt]{article}
\usepackage[margin=1in]{geometry}
\usepackage{graphicx}
\usepackage{booktabs}
\usepackage{amsmath}
\usepackage{hyperref}
\usepackage{natbib}

\title{The Null Is Part of the Instrument: Construction-Matched Controls in AI-Assisted Experimental Research}

\author{Physis Research Program}
\date{\today}

\begin{document}
\maketitle

\begin{abstract}
In AI-assisted empirical research, the construction-matched null is not a
formality but the primary instrument. We document twelve cases from a
three-month research program where the null's design determined the
verdict: a positive result survived only when a control specifically
capable of killing the claimed mechanism failed to kill it. We propose
that null design is a first-class research concern and that the quality
of a study is determined by the quality of its nulls, not by the
sophistication of its models.
\end{abstract}

\section{Introduction}

The standing rule of experimental software research is: a positive result
is an artifact until a control specifically capable of killing the
claimed mechanism fails to kill it. This rule is easy to state and hard
to practice. In AI-assisted research, where agents design, implement, and
interpret their own experiments, the temptation to run a token control
--- one that cannot fail --- is strong.

We document twelve cases where the null's design determined the verdict.
In six cases, a pre-registered null killed a positive result. In four
cases, a control that appeared valid was shown to be invalid by a second
test. In two cases, a control that was valid at one model scale became
invalid at another.

\section{The Null Is the Instrument}

Consider the canonical experiment: a system claims to detect misfiled
ontology entries at AUC 0.880. The claim is tested against a null. If
the null is construction-matched --- if it places errors where the system
cannot detect them --- then the 0.880 is meaningful. If the null is
token --- if it places errors at random, where any distance-based scorer
can detect them --- then the 0.880 is a measurement of the null's
weakness, not the system's strength.

In our research, the injection policy (random vs.\ nearest-cell) changed
the AUC from 0.880 to 0.506. The system was the same. The metric was the
same. Only the null changed. The null is the instrument.

This is not a subtle point. In the perspective-discovery track, the
central negative result --- cross-embedder ontology agreement is near
chance --- was established not by a new clustering algorithm but by a new
null: conditioning on the stronger support instead of the baseline. The
mechanism (D = |s1 - s2| is a folded read-out of s2) was identified by
the null, not by the positive result.

\section{Twelve Cases}

We organize the twelve cases by the phase of research they affect.

\subsection{Design Phase}

\textbf{Case 1: Injection distribution (E14).} Claim: geometry scores
misfilings at AUC 0.880. Null: injection at nearest alternative cell.
Result: AUC 0.506. The null killed the claim by changing only where
errors were placed.

\textbf{Case 2: Fixture circularity (E83).} Claim: antonym rule detects
contradiction. Null: hand-written fixtures without lexical antonyms.
Result: recall 0.143 vs.\ 0.998 on generated fixtures. The null killed
the claim by removing the generator's circularity.

\textbf{Case 3: Adversarial fixtures (E85).} Claim: structured rules
achieve precision 1.000. Null: adversarial fixtures designed to fire
each rule wrongly. Result: precision 0.000. The null killed the claim by
inviting the rules to fail.

\textbf{Case 4: Scale-dependent controls (E19b).} Claim: five congruent
domains improve abstraction. Null: control arm B5c (five coupled domains,
each failing a different clause). Result: the control fires at 4B but is
valid at 9B. The null's validity was a property of the stimulus and the
model together.

\subsection{Implementation Phase}

\textbf{Case 5: Harness parsing (E74).} Claim: Cline is a reliable
agent. Null: check whether shell actions have output. Result: failure
rate 0.005 (parser found no output) vs.\ 0.070 (parser reads output).
The null killed the claim by checking the parser, not the agent.

\textbf{Case 6: Duplicate sessions (E73).} Claim: the corpus has 402
sessions. Null: check for byte-identical session IDs across machines.
Result: 188 duplicate records (47\% inflation). The null killed the
count by checking for duplicates, not by trusting the file count.

\textbf{Case 7: Truncated embeddings (E78).} Claim: revisions form a
tight region. Null: check whether the embedding covers the whole file.
Result: first 1200 bytes produced identical vectors (within-object
distance 0.0000). The null killed the claim by measuring the embedding
coverage.

\textbf{Case 8: Parser measuring the wrong field (E76).} Claim: coherence
improves after perturbation. Null: check whether the parser reads
coherence or node count. Result: parser picked up node count. The null
killed the claim by inspecting the parser.

\subsection{Interpretation Phase}

\textbf{Case 9: Confound check that stops early (E17).} Claim: ANCHORVOC
is not circular. Null: three tests (predicts rater's choice; survives
half the hints; same on hint-leaking entries). Result: Test 3 fails.
The null killed the claim by requiring all three tests to pass.

\textbf{Case 10: Baseline drift (E87).} Claim: pooled baseline scores
0.400. Null: rerun the baseline in the same script. Result: 0.200. The
null killed the claim by rerunning, not by trusting the previous number.

\textbf{Case 11: Aggregation hiding components (E81).} Claim: static
rules find no structure (ratio 0.96). Null: decompose by cue arity.
Result: multi-word cues 15:1, single-word cues 0.92:1. The null killed
the claim by decomposing the aggregate.

\textbf{Case 12: Pre-flight ignored (E10).} Claim: the experiment is
powered. Null: power analysis. Result: STOP. The run proceeded anyway.
The null was overridden by its author --- a control the author ignores is
not a control.

\section{Principles}

From these cases we derive three principles:

\textbf{P1: The null must be construction-matched.} A null that places
errors where the system cannot detect them is not a control; it is a
tautology. The null must perturb the claimed mechanism.

\textbf{P2: The null's validity must be tested.} A control that passes
one test and fails another is not a valid control. The workflow must
require all pre-registered tests to pass, not just the first reassuring
one.

\textbf{P3: Baselines must be rerun in the same execution.} A baseline
number cited from a previous run is an assertion about the past, not a
measurement. Drift of 0.2 between runs is not rare; it is the norm.

\section{Related Work}

The importance of separating exploratory from confirmatory analysis is
long-established in statistics \citep{tukey1980we}, and the cost of
undisclosed analytic flexibility is well documented in psychology
\citep{simmons2011false}. The novel contribution here
is the \textbf{executable} null: a control that is run in the same
script, on the same data, with the same code path, and whose result
determines the verdict. In AI-assisted research, where agents can
reinterpret results after the fact, the executable null is the only kind
that cannot be argued with.

The construction-matched null is related to the adversarial input used to
probe machine-learning systems: both are designed to expose the limits of
a system rather than to confirm its strengths. The difference is that the
construction-matched null is pre-registered: it is specified before the
experiment is run, not invented after the result is known. We cite no
specific work for the adversarial-input analogy because we have not
verified a primary source for it in this revision.

\section{Limitations}

The twelve cases are from one project, one user, one machine fleet. The
frequency of each failure mode in other contexts is unknown. The
principles are derived from cases where the null caught the error; we do
not have a systematic study of cases where the null failed to catch an
error that a human would have caught.

\section{Conclusion}

The construction-matched null is the primary instrument in AI-assisted
empirical research. Its design determines what can be concluded. A
positive result is an artifact until a null specifically capable of
killing the claimed mechanism fails to kill it. Twelve cases demonstrate
that this is not a philosophical preference but a practical necessity.

\begin{thebibliography}{99}

\bibitem[Tukey(1980)]{tukey1980we}
Tukey, J.~W. (1980).
\newblock We need both exploratory and confirmatory.
\newblock \emph{The American Statistician}, 34(1), 23--25.
\newblock DOI: 10.1080/00031305.1980.10482706.

\bibitem[Simmons et al.(2011)]{simmons2011false}
Simmons, J.~P., Nelson, L.~D., \& Simonsohn, U. (2011).
\newblock False-positive psychology: Undisclosed flexibility in data
collection and analysis allows presenting anything as significant.
\newblock \emph{Psychological Science}, 22(11), 1359--1366.
\newblock DOI: 10.1177/0956797611417632.

\end{thebibliography}

\end{document}
